% QUARANTINED RESEARCH ONLY -- 29 August 2026.
% Revised6 standalone repair module. Not a live-book input and not production-verified.

\documentclass[11pt]{article}
\usepackage[T1]{fontenc}
\usepackage{lmodern}
\usepackage{amsmath,amssymb,amsthm,mathtools}
\usepackage[a4paper,margin=27mm]{geometry}
\usepackage{enumitem}
\usepackage{needspace}
\usepackage[hidelinks,hypertexnames=false]{hyperref}
\usepackage[nameinlink,capitalize,noabbrev]{cleveref}

\newtheorem{theorem}{Theorem}
\newtheorem{lemma}[theorem]{Lemma}
\newtheorem{corollary}[theorem]{Corollary}
\newtheorem{remark}[theorem]{Remark}

\newcommand{\MA}{\operatorname{MA}}
\newcommand{\ph}{\operatorname{ph}}
\newcommand{\cM}{\mathcal M}
\newcommand{\cR}{\mathcal R}
\newcommand{\R}{\mathbf R}
\newcommand{\C}{\mathbf C}
\newcommand{\eps}{\varepsilon}
\newcommand{\norm}[1]{\lVert#1\rVert}
\newcommand{\abs}[1]{\lvert#1\rvert}
\setlength{\emergencystretch}{2em}

\title{P17.1 repair module, revised6:\\
Powered Archimedean normalization, Philippon--Gauss endpoint,\\
and uniform CLN passage}
\author{}
\date{29 August 2026\\\small Quarantined research only; no live integration or production submission}

\begin{document}
\maketitle

\begin{abstract}
This module repairs the powered Archimedean and uniform-density part of the
proposed P17.1 bridge.  It derives the metric unpowering convention from the
explicit fibre formula in Pille--Schneider Proposition~3.2, connects that
formula to their Definition~4.8 and to Boucksom--Eriksson's usual
Archimedean normalization, and only then proves the Monge--Ampere and Deligne
scaling laws.  It also imports the separately staged weighted
Philippon--Gauss theorem and reconstructs its non-Archimedean endpoint,
proves a two-sided CLN bound uniform in the power on the exact semipositive
slot-zero domain required by the cited inequalities, and closes the density passage by a
coupled cofinal product net.  The module remains conditional on the cited
standard inputs and on finite-level continuity.  It neither constructs nor
compares the relative resultant tensor and therefore does not certify the
whole P17.1 bridge.
\end{abstract}

\section{The power parameter and the exact fibre convention}

Let $A$ be an integral Banach ring.  For every $x\in\cM(A)$, set
\begin{equation}
 \rho(x):=\frac{\log\max\{1,\abs{2}_x\}}{\log 2}.
\tag{1.1}\label{eq:rho-base}
\end{equation}
This is a continuous function on $\cM(A)$.

\begin{lemma}
\label{lem:rho-split}
One has $\rho(x)=0$ exactly when the absolute value on the completed residue
field $\mathcal H(x)$ is non-Archimedean.  If $\rho(x)>0$, then
$0<\rho(x)\leq1$, and there is a character into $\C$, unique up to complex
conjugation, such that after completing the fraction field
\begin{equation}
 \abs{a}_x=\abs{\chi_x(a)}_\infty^{\rho(x)}.
\tag{1.2}\label{eq:character-power}
\end{equation}
Thus $\mathcal H(x)$ is $\R$ or $\C$, with the
$\rho(x)$-th power of the usual absolute value.
\end{lemma}

\begin{proof}
If $\abs{2}_x\leq1$, binary expansion first gives
$\abs{n}_x\leq1$ for all positive integers $n$; applying the binomial
formula to $(a+b)^m$ and taking $m$-th roots gives the ultrametric
inequality.  The converse follows from
$\abs{2}_x\leq\max\{\abs1_x,\abs1_x\}=1$.

If $\abs2_x>1$, the base-expansion proof of Ostrowski's theorem gives
$\abs n_x=n^{\rho(x)}$.  The ordinary triangle inequality gives
$\abs2_x\leq2$, hence $0<\rho(x)\leq1$.  The unpowered function
$v(a)=\abs a_x^{1/\rho(x)}$ is multiplicative, is usual on $\mathbf Q$,
and satisfies the triangle inequality: expanding $(a+b)^m$ yields
\[
 \abs{a+b}_x^m\leq(m+1)
 \bigl(v(a)+v(b)\bigr)^{m\rho(x)},
\]
and the power $1/(m\rho(x))$, followed by $m\to\infty$, proves the
claim.  The completion is a commutative real Banach division algebra, hence
$\R$ or $\C$ by Gelfand--Mazur.  The only isometric ambiguity in the
complex case is conjugation.  This is the argument in the separately staged
proof imported for \cref{thm:weighted-PG} below.
\end{proof}

Fix henceforth an Archimedean point and abbreviate
$K_\rho=(K,\abs{\cdot}_\infty^\rho)$, where
$K\in\{\R,\C\}$ and $0<\rho\leq1$.  A metric weight is
additive: for a local frame $e$,
\[
 \phi_\rho(e):=-\log\norm e_\rho,
 \qquad
 \phi_\rho(ue)=\phi_\rho(e)-\log\abs u_\rho.
\tag{1.3}\label{eq:metric-transition}
\]
Its unpowered metric is
\begin{equation}
 \overline\phi(e):=\rho^{-1}\phi_\rho(e),
 \qquad
 \norm e_{\overline\phi,\infty}:=\norm e_{\phi_\rho,\rho}^{1/\rho}.
\tag{1.4}\label{eq:metric-unpower}
\end{equation}
Division of~\eqref{eq:metric-transition} by $\rho$ gives the usual metric
transition law.

\begin{lemma}
\label{lem:PS-transport}
Under the canonical homeomorphism from a usual complex or real
analytification to the $K_\rho$-analytification, Pille--Schneider's
restricted fibre metric is transported to the usual complex metric by
\eqref{eq:metric-unpower}.  Their fibrewise Monge--Ampere measure in
Definition~4.8 is therefore Boucksom--Eriksson's Archimedean measure of the
unpowered metrics.  This is a consequence of their stated fibre convention,
not a new formal definition of a symbol $dd^c_\rho$.
\end{lemma}

\begin{proof}
The exponent is explicit in Pille--Schneider Proposition~3.2.  In their
hybrid-disc notation, put
\[
 \rho_t:=\frac{\log r}{\log\abs t}\quad(0<\abs t\leq r<1).
\]
For the canonical homeomorphism $\beta_t:X_t\to X_t^{\rm hyb}$, their
proof gives, for every local regular function $f$,
\begin{equation}
 \abs{f(\beta_t(z))}=\abs{f(z)}_\infty^{\rho_t},
 \qquad
 \norm{s(z)}_{\phi_t}
 =\norm{s(\beta_t(z))}_{\phi}^{1/\rho_t}.
\tag{1.5}\label{eq:PS32-transport}
\end{equation}
Equivalently,
$\log\norm{s}_{\phi_t}/\log_r\abs t$
extends as the hybrid logarithmic norm.  Thus the raw hybrid fibre weight is
$\rho_t$ times the ordinary complex weight, exactly as in
\eqref{eq:metric-unpower}.

For an arbitrary Archimedean $x\in\cM(A)$, the character
\eqref{eq:character-power} gives the same canonical map: on every affine
chart it sends a complex point $z$ to
$f\mapsto\abs{\chi_x(f)(z)}_\infty^{\rho(x)}$.  Formula
\eqref{eq:metric-unpower} is therefore the unique metric transport compatible
with local frames and this seminorm map; a different representing character
only conjugates the construction.

This transport also preserves the psh class actually used by
Pille--Schneider.  For a tropical Fubini--Study metric,
\[
 \rho^{-1}\frac1m\max_j
 \{\log\abs{s_j}_\rho+a_j\}
 =\frac1m\max_j
 \{\log\abs{s_j}_\infty+a_j/\rho\},
\tag{1.6}\label{eq:FS-transport}
\]
which is usual complex psh by their Theorem~2.11.  Their
Proposition~2.21 and multiplication by the fixed scalar $1/\rho$ transport
uniform limits as well.

Finally, Pille--Schneider Definition~4.8 defines the family measure at $x$
by restricting the metrics to $X_{\mathcal H(x)}^{\rm an}$ and applying the
valued-field Monge--Ampere operator of their Section~4.1.  On an
Archimedean field Boucksom--Eriksson first identifies the field with
$\R$ or $\C$ only \emph{up to normalization of its absolute value}; their
Archimedean current is the classical Bedford--Taylor current on the usual
complex space.  Equations~\eqref{eq:PS32-transport}--\eqref{eq:FS-transport}
identify the metric supplied to that current as $\overline\phi$.

Two further source tests make this reading fail-closed.  First,
Pille--Schneider Theorem~4.11 writes the Archimedean members of a hybrid
family as $(dd^c\phi_t)^n$, where Proposition~3.2 has already defined
$\phi_t=\rho_t^{-1}\phi|_{X_t^{\rm hyb}}$.  Thus their displayed complex
operator acts on the unpowered metric.

Second, Boucksom--Eriksson Theorem~8.14 is stated over an arbitrary complete
valued field, and its proof uses
\[
 dd^c\log\norm{s}_{\phi}\wedge\delta_X
 =\delta_Z-dd^c\phi\wedge\delta_X.
\tag{1.7}\label{eq:PL-source}
\]
For $K_\rho$, both
$\log\norm{s}_{\phi_\rho,\rho}$ and $\phi_\rho$ are $\rho$ times their
unpowered counterparts.  Applying the raw Euclidean operator to those
powered weights would turn the left side into
$\rho\delta_Z-\rho dd^c_\infty\overline\phi$, while the right side demanded
by~\eqref{eq:PL-source} would be
$\delta_Z-\rho dd^c_\infty\overline\phi$.  These differ for
$\rho\ne1$ and a nonzero divisor.  The source identity is restored exactly
by the transported operator
\begin{equation}
 d^c_\rho\phi_\rho:=d^c_\infty(\rho^{-1}\phi_\rho),
 \qquad
 dd^c_\rho\phi_\rho:=dd^c_\infty(\rho^{-1}\phi_\rho).
\tag{1.8}\label{eq:ddc-unpacking}
\end{equation}
Equation~\eqref{eq:ddc-unpacking} is our explicit unpacking of the source
convention, not a formula printed verbatim in either source.
Boucksom--Eriksson Proposition~8.3(ii), which fixes total mass to the
algebraic intersection number, supplies the independent mass check: the raw
Euclidean operator on $\rho\overline\phi$ would multiply an $n$-fold mass
by $\rho^n$.  Hence Definition~4.8 has the transported reading asserted in
the lemma.
\end{proof}

\section{PN1, dimension zero, and recursive PN2}

Let $X/K_\rho$ be projective of dimension $n$; nilpotents, generic
multiplicities, and several top-dimensional components are allowed.

\begin{theorem}
\label{thm:powered-normalization}
For continuous psh-regularizable metrics $\phi_{i,\rho}$, put
$\overline\phi_i=\rho^{-1}\phi_{i,\rho}$.  Under the canonical
identification of analytic spaces,
\begin{equation}
 \MA_\rho(\phi_{1,\rho},\ldots,\phi_{n,\rho})
 =\MA_\infty(\overline\phi_1,\ldots,\overline\phi_n).
\tag{PN1}\label{eq:PN1-r3}
\end{equation}
All generic lengths in the fundamental cycle occur on both sides.  If $e$
is a nonzero vector in the Deligne line $I_X(L_0,\ldots,L_n)$, then
\begin{equation}
 \lambda_\rho(e;\phi_{0,\rho},\ldots,\phi_{n,\rho})
 =\rho\lambda_\infty
 (e;\overline\phi_0,\ldots,\overline\phi_n).
\tag{PN2}\label{eq:PN2-r3}
\end{equation}
Equivalently, the recursively defined local height has exactly one factor
$\rho$ in every dimension.
\end{theorem}

\begin{proof}[Proof of PN1]
\cref{lem:PS-transport} identifies the smooth complex fibre measure with
\[
 dd^c_\infty\overline\phi_1\wedge\cdots\wedge
 dd^c_\infty\overline\phi_n.
\]
For a general scheme, Boucksom--Eriksson use the fundamental cycle
\[
 [X]_n=\sum_Zm_Z[Z],
 \qquad
 m_Z=\operatorname{length}(\mathcal O_{X,\eta_Z}).
\]
Their measure is the sum of the component measures multiplied by $m_Z$.
Applying the transport on each reduced top-dimensional component proves the
smooth statement with all generic lengths.

For continuous psh-regularizable metrics, uniformly approximate by smooth
semipositive nets.  Division by the fixed $\rho$ preserves uniform
convergence.  Boucksom--Eriksson Theorem~8.4, used in Pille--Schneider
Theorem~4.3, gives the unique weakly continuous extension and hence
\eqref{eq:PN1-r3}.  Over $\R$, both sources define the measure as the
push-forward from complexification.  Complexification commutes with the
transport and introduces no factor two.
\end{proof}

\begin{lemma}
\label{lem:zero-anchor-r3}
Let $X=\operatorname{Spec}A$ be finite and flat over $K_\rho$, and let
$\phi_{0,\rho}$ be a continuous dpsh metric on $L_0$.  Put
$\overline\phi_0=\rho^{-1}\phi_{0,\rho}$.  Under the algebraic
identification of the Deligne line with the norm line,
\begin{equation}
 \log\norm{\langle s\rangle}_{\langle\phi_{0,\rho}\rangle,\rho}
 =\rho\,
 \log\norm{\langle s\rangle}_{\langle\overline\phi_0\rangle,\infty}
\tag{Z0}\label{eq:zero-metric-r3}
\end{equation}
for every trivializing section $s$ of $L_0$.  Consequently the additive
Deligne weight and every dimension-zero local height are multiplied by
exactly $\rho$, with all residue degrees and nonreduced Artin lengths.
For $K=\R$, the same assertion uses the real-fibre push-forward convention
and has no factor two.
\end{lemma}

\begin{proof}
The algebraic input is Boucksom--Eriksson Example~8.15: in dimension zero
$\langle L_0\rangle=N_{A/K}(L_0)$ and
$\langle us\rangle=N_{A/K}(u)\langle s\rangle$.  Ducrot
Proposition~6.1.1(3) fixes the positive algebraic norm orientation, but
neither algebraic statement fixes the multiplicative constant of the
metric.

That constant is fixed in the proof of Boucksom--Eriksson Theorem~8.16.
Their formula~(8.11), reconstructed in the present norm notation, is
\begin{equation}
 \log\norm{\langle s\rangle}_{\langle\phi_0\rangle}
 =\int_{X^{\rm an}}\log\norm{s}_{\phi_0}\,\delta_X.
\tag{8.11$'$}\label{eq:BE811-r3}
\end{equation}
Write $A=\prod_jA_j$, with $A_j$ local, residue field $K_j$, and Artin
length $m_j$.  Boucksom--Eriksson Example~8.2 gives, on the corresponding
analytic points $x_j$,
\begin{equation}
 \delta_X=\sum_j(\dim_KA_j)\delta_{x_j}
 =\sum_jm_j[K_j:K]\,\delta_{x_j}.
\tag{2.2}\label{eq:zero-cycle-r3}
\end{equation}
Under the canonical identification of the powered and unpowered analytic
spaces, this measure is unchanged, while metric unpowering gives
$\log\norm{s(x_j)}_{\phi_{0,\rho},\rho}
 =\rho\log\norm{s(x_j)}_{\overline\phi_0,\infty}$.
Applying~\eqref{eq:BE811-r3} twice and then~\eqref{eq:zero-cycle-r3}
therefore yields
\begin{align*}
 \log\norm{\langle s\rangle}_{\langle\phi_{0,\rho}\rangle,\rho}
 &=\sum_jm_j[K_j:K]\,
   \log\norm{s(x_j)}_{\phi_{0,\rho},\rho}\\
 &=\rho\sum_jm_j[K_j:K]\,
   \log\norm{s(x_j)}_{\overline\phi_0,\infty}\\
 &=\rho\,
   \log\norm{\langle s\rangle}_{\langle\overline\phi_0\rangle,\infty}.
\end{align*}
This proves~\eqref{eq:zero-metric-r3} and visibly preserves both the residue
degree $[K_j:K]$ and the Artin length $m_j$.  Boucksom--Eriksson verify the
frame compatibility by their formula~(8.12),
$\int\log\abs u\,\delta_X=\log\abs{N_{A/K}(u)}$; after powering, both
sides acquire the same factor $\rho$.  Finally, their Remark~8.1 defines
the real measure as the push-forward from complexification.  Thus the
degree already present in~\eqref{eq:zero-cycle-r3} is retained and no
additional factor two appears.
\end{proof}

\medskip
\begin{proof}[Proof of PN2]
Define the candidate powered Deligne weight by
\[
 \widetilde\lambda_\rho
 (e;\phi_{0,\rho},\ldots,\phi_{n,\rho})
 :=\rho\lambda_\infty
 (e;\overline\phi_0,\ldots,\overline\phi_n).
\tag{2.1}\label{eq:PN2-candidate-r3}
\]
If $e'=ue$, then
$
\widetilde\lambda_\rho(e')=\widetilde\lambda_\rho(e)
-\rho\log\abs u_\infty
=\widetilde\lambda_\rho(e)-\log\abs u_\rho
$, so this is a metric for the powered absolute value.
At fixed $\rho$, unpowering is additive under tensor products; hence
symmetry and multilinearity of the usual Deligne metric give the same
properties for $\widetilde\lambda_\rho$.

Let $s\in H^0(X,L_0)$ be regular, with divisor $Z$.  The usual additive
restriction formula is
\begin{align*}
 \lambda_\infty(e;\overline\phi_0,\ldots,\overline\phi_n)
 ={}&\lambda_\infty(e|_Z;
       \overline\phi_1|_Z,\ldots,\overline\phi_n|_Z)\\
 &-\int_X\log\norm{s}_{\overline\phi_0,\infty}
       \,\MA_\infty(\overline\phi_1,\ldots,\overline\phi_n).
\end{align*}
Multiplication by $\rho$, the identity
$
 \log\norm{s}_{\phi_{0,\rho},\rho}
 =\rho\log\norm{s}_{\overline\phi_0,\infty}
$, and PN1 turn this into the powered restriction formula.  The
dimension-zero lemma supplies the base case, and the uniqueness clause in
Boucksom--Eriksson Theorem~8.16 identifies the candidate with the target
Deligne metric.

The same calculation is the local-height induction: the restricted height
has one factor $\rho$, while the integral has one factor from its Green
function and none from its measure.  No new factor is accumulated at the
next dimension.
\end{proof}

\section{The weighted Philippon--Gauss theorem, including the endpoint}

The next theorem is imported from the separately staged, independently
verified blueprint
\path{ch66_weighted_philippon_gauss_statement.tex}, with its full proof in
\path{ch66_weighted_philippon_gauss_proof.tex}.  The conclusion is restated
here, and the load-bearing character and endpoint arguments are reconstructed
below.  This reconstruction is not asserted to be an exact-byte copy of the
staged full proof.

\medskip
\Needspace{0.32\textheight}
\begin{theorem}
\label{thm:weighted-PG}
Let $A$ be a Banach ring and $U\subset\cM(A)$ open.  Fix finitely many
blocks
\[
 \mathbf U_i=(U_{i0},\ldots,U_{i,N_i-1}),\qquad N_i\geq1,
\]
nonnegative block degrees $d_i$, a finite set $S$ of multiexponents of
block degree $(d_i)_i$, and real weights $w_\alpha$.  Let
$c_\alpha\in\mathcal O(U)$, assume the coefficient tuple is nowhere zero,
and assume locally that all $c_\alpha$ belong to one Banach algebra of
analytic functions on a Banach analytic domain.

Use $\rho(x)$ from~\eqref{eq:rho-base}.  If $\rho(x)=0$, set
\begin{equation}
 \Phi(x):=\max_{\alpha\in S}
 \bigl(\log\abs{c_\alpha(x)}_x+w_\alpha\bigr),
 \qquad \log0:=-\infty.
\tag{PG$_0$}\label{eq:PG-zero}
\end{equation}
If $\rho(x)>0$, choose a local character satisfying
\eqref{eq:character-power}, put
\begin{equation}
 F_x:=\sum_{\alpha\in S}\chi_x(c_\alpha)
       \exp(w_\alpha/\rho(x))\mathbf U^\alpha,
 \qquad
 \Phi(x):=\rho(x)\ph(F_x),
\tag{PG$_+$}\label{eq:PG-positive}
\end{equation}
where $\ph$ is Philippon's logarithmic sphere measure, including its
harmonic correction.  Then $\Phi$ is choice-independent and continuous on
$U$.

For a nowhere-vanishing analytic function $\lambda$,
\begin{equation}
 \Phi(\lambda F)=\log\abs\lambda+\Phi(F).
\tag{3.1}\label{eq:PG-homogeneity}
\end{equation}
Consequently, if $\mathcal L$ is an analytic line bundle and
$R=eF_e$ is an $\mathcal L$-valued polynomial, then for $e'=ue$
\begin{equation}
 F_{e'}=u^{-1}F_e,
 \qquad
 \Phi_{e'}=\Phi_e-\log\abs u,
\tag{3.2}\label{eq:PG-frame}
\end{equation}
and these weights define a continuous additive metric on $\mathcal L$.
\end{theorem}

\begin{proof}
\cref{lem:rho-split} gives the two loci and the character, unique up to
conjugation.  To control those characters in families, work locally in the
finitely generated $\mathbf Z$-algebra $B$ generated by the coefficients
and the finitely many localization elements.  If $x_\lambda\to x$ and
$\rho(x)>0$, then eventually $\rho(x_\lambda)\geq\rho(x)/2$.  The values of
chosen characters on generators of $B$ lie in one compact polydisc.  Every
subnet therefore has a further subnet converging pointwise on $B$ to a
character representing $x$; uniqueness makes that limit $\chi_x$ or its
complex conjugate.

On every punctured fixed-type coefficient space,
$\ph$ is finite, continuous, log-homogeneous, and invariant under
conjugation.  The coefficient continuity follows from Zhang's continuity of
the Chow metric; padding the blocks and multiplying by fixed coordinate
monomials reduces unequal block sizes and degrees to equal Chow type.  Hence
\eqref{eq:PG-positive} is independent of the representing character and is
continuous along each of the further subnets above.  The subnet criterion for
convergence then proves continuity on the open locus $\rho>0$.

It remains to prove the endpoint rather than infer it by division by
$\rho$.  D'Andrea--Krick--Sombra give, for the fixed block type,
\begin{equation}
 \left|\ph(F)-\log\norm F_{\max}\right|\leq
 B,\qquad
 B:=2\sum_i\log(N_i+1)d_i.
\tag{3.3}\label{eq:PG-coefficient-bound}
\end{equation}
Define the finite continuous function
\[
 M(x):=\max_{\alpha\in S}
 \bigl(\log\abs{c_\alpha(x)}_x+w_\alpha\bigr).
\]
It is finite because the coefficient tuple is nowhere zero.  On
$\rho(x)>0$, there is the exact identity
\begin{align}
 \rho(x)\log\norm{F_x}_{\max}
 &=\rho(x)\max_\alpha\left(
 \log\abs{\chi_x(c_\alpha)}_\infty+
 \frac{w_\alpha}{\rho(x)}\right)\notag\\
 &=M(x).
\tag{3.4}\label{eq:PG-exact-max}
\end{align}
Multiplying~\eqref{eq:PG-coefficient-bound} by $\rho(x)$ gives
\begin{equation}
 \abs{\Phi(x)-M(x)}\leq B\rho(x).
\tag{3.5}\label{eq:PG-boundary}
\end{equation}
At $\rho=0$, definition~\eqref{eq:PG-zero} is exactly $\Phi=M$.  Thus for
any net $x_\lambda\to x$ with $\rho(x)=0$, continuity of $M$ and
$\rho$ and estimate~\eqref{eq:PG-boundary} give
$\Phi(x_\lambda)\to M(x)=\Phi(x)$, whether the net approaches through
either locus.  This proves the endpoint and global continuity.

Log-homogeneity of $\ph$ gives~\eqref{eq:PG-homogeneity} on
$\rho>0$, while the maximum formula gives it at $\rho=0$.  Finally
$eF_e=e'F_{e'}=ueF_{e'}$, which proves~\eqref{eq:PG-frame} with the
stated sign.
\end{proof}

\begin{corollary}
\label{cor:powered-PG}
Let $P$ be a nonzero multihomogeneous resultant polynomial in blocks of
dimensions $r_i+1$ and degrees $\delta_i$, in coordinates orthonormal for
the unpowered Hermitian norms.  At $\rho>0$,
\begin{equation}
 \rho\ph(P)=\rho\int_{\prod_iS^{2r_i+1}}
 \log\abs{P(z)}\,\mathrm d\eta(z)
 +\frac\rho2\sum_i\delta_i\sum_{\ell=1}^{r_i}\frac1\ell.
\tag{PH$_\rho$}\label{eq:powered-PH-r3}
\end{equation}
If $\cR=e^{\otimes Q}F_e$ and $e'=ue$, then on all of $\cM(A)$
\begin{equation}
 F_{e'}=u^{-Q}F_e,
 \qquad
 \Phi_{e'}=\Phi_e-Q\log\abs u.
\tag{FL}\label{eq:frame-law-r3}
\end{equation}
\smallskip
At $\rho=0$, $\Phi$ is the weighted Gauss maximum
\eqref{eq:PG-zero}.
\end{corollary}

\begin{proof}
Chen--Moriwaki Theorem~3.8.7 gives the displayed sphere integral plus harmonic
correction for the usual absolute value.  PN2 multiplies the entire local
height by $\rho$, including the correction.  If an unpowered basis vector
has norm $\exp(-a_{i,\beta}/\rho)$, its orthonormalization multiplies it by
$\exp(a_{i,\beta}/\rho)$; hence a coefficient of multiexponent $\alpha$
gets the exact factor
\[
 \exp(w_\alpha/\rho),\qquad
 w_\alpha=\sum_{i,\beta}\alpha_{i,\beta}a_{i,\beta}.
\]
For real fibres, complexification and conjugation invariance give the same
formula.  Apply \cref{thm:weighted-PG} to the line bundle
$\mathcal L^{\otimes Q}$, whose frame changes by $u^Q$; this proves
\eqref{eq:frame-law-r3} and supplies the endpoint.
\end{proof}

\section{Two-sided powered CLN and family uniformity}

For every $0\leq i\leq n$, let
$\phi_{i,\rho},\phi'_{i,\rho}$ be continuous psh-regularizable
semipositive metrics on $L_i$, and put
$f_\rho=\phi_{0,\rho}-\phi'_{0,\rho}$.  In particular, the slot-zero
test is a difference of two metrics in the semipositive domain of the CLN
inequality; it is not an arbitrary continuous function.  Define
\begin{equation}
 C:=
 \begin{cases}
 0,&n=0,\\[2mm]
 2\displaystyle\max_{1\leq i\leq n}
 (L_0L_1\cdots\widehat{L_i}\cdots L_n),&n\geq1.
 \end{cases}
\tag{4.1}\label{eq:CLN-C-r3}
\end{equation}

\begin{lemma}
\label{lem:CLN-r3}
Assume explicitly that both metrics in every pair
$\phi_{i,\rho},\phi'_{i,\rho}$, including $i=0$, are continuous
psh-regularizable and semipositive.  For $0<\rho\leq1$,
\begin{align}
 &\left|\int_Xf_\rho\,
 \MA_\rho(\phi_{1,\rho},\ldots,\phi_{n,\rho})
 -\int_Xf_\rho\,
 \MA_\rho(\phi'_{1,\rho},\ldots,\phi'_{n,\rho})\right|\notag\\
 &\hspace{30mm}\leq
 C\sum_{i=1}^n
 \norm{\phi_{i,\rho}-\phi'_{i,\rho}}_\infty.
\tag{CLN$_\rho$}\label{eq:CLN-r3}
\end{align}
The constant is independent of $\rho$.
\end{lemma}

\begin{proof}
If $n=0$, there are no metric slots to change, the two integrals coincide,
and both sides are zero by definition~\eqref{eq:CLN-C-r3}.

Assume $n\geq1$ and first suppose all metrics are smooth semipositive.
Put $\overline\phi_i=\rho^{-1}\phi_{i,\rho}$ and
$\overline\phi'_i=\rho^{-1}\phi'_{i,\rho}$ for every $0\leq i\leq n$,
and $\overline f=\overline\phi_0-\overline\phi'_0$.
The convention transport preserves semipositivity, and PN1 makes the left
side equal to
\[
 \rho\left|\int_X\overline f\,
 \MA_\infty(\overline\phi_1,\ldots)
 -\int_X\overline f\,
 \MA_\infty(\overline\phi'_1,\ldots)\right|.
\]
Boucksom--Eriksson Lemma~8.6 applies because the slot-zero pair
$(\overline\phi_0,\overline\phi'_0)$ is also semipositive.  It gives one
inequality with constant~\eqref{eq:CLN-C-r3}.  Keep that test pair fixed and
interchange primed and unprimed metrics only in slots $1,\ldots,n$ to get
the opposite inequality.  Therefore
\begin{align*}
 \rho\left|\int\overline f\,\MA_\infty(\overline\phi_i)
 -\int\overline f\,\MA_\infty(\overline\phi'_i)\right|
 &\leq\rho C\sum_i
 \norm{\overline\phi_i-\overline\phi'_i}_\infty\\
 &=C\sum_i
 \norm{\phi_{i,\rho}-\phi'_{i,\rho}}_\infty.
\end{align*}
This is the two-sided cancellation: the test function contributes $\rho$,
and unpowering a metric difference contributes $1/\rho$.

For the stated continuous domain, uniformly approximate all
$2(n+1)$ metrics, including the slot-zero pair, by smooth semipositive
metrics.  Boucksom--Eriksson Theorem~8.4 gives weak convergence of the
Monge--Ampere measures, while uniform convergence of the slot-zero
differences passes the integrals to the limit.  The same constant $C$
therefore survives.  This argument supplies no estimate for an arbitrary
continuous test function outside a difference of two such semipositive
metrics.
\end{proof}

\medskip
\Needspace{0.22\textheight}
\begin{corollary}
\label{cor:uniform-family-r3}
Let $X/A$ be flat and projective, let the line bundles above come from $X$,
and assume that for every $0\leq i\leq n$ both metrics in the pair
$\phi_i,\phi'_i$ are continuous psh-regularizable and semipositive on the
fibres.  The same constant works on every fibre for the test
$\phi_0-\phi'_0$.  It is not asserted for an arbitrary continuous test
function.  At $\rho(x)>0$ this is \cref{lem:CLN-r3}; at
$\rho(x)=0$ it is Pille--Schneider Lemma~4.10 and its continuous
semipositive extension.  No division by $\rho$ occurs on the endpoint
locus.
\end{corollary}

\begin{proof}
For $n=0$, $C=0$.  For $n\geq1$, the constant is made only from mixed
algebraic intersection numbers, which are fibre-independent by flatness.
\cref{lem:CLN-r3} has no residual power parameter.  On the
non-Archimedean locus, Pille--Schneider Lemma~4.10 starts with FS metrics in
every slot, including the slot-zero pair.  Uniform FS approximation of all
these semipositive metrics and the same inequality pass the estimate to the
stated continuous domain.
\end{proof}

\section{Uniform density by a coupled cofinal product net}

Assume the finite-level fibre integral is continuous for every DFS test
function and every tuple of tropical FS metrics.  Let
$M=(L_1\cdots L_n)\geq0$ be the common fibre mass.  Choose a net of DFS
functions $\psi_q\to\psi$ uniformly.  By the definition of DFS, after
transporting one line bundle through the isomorphism in the DFS presentation,
write
\[
 \psi_q=\alpha_q-\beta_q
\]
with $\alpha_q,\beta_q$ tropical FS metrics on one semiample line bundle
$H_q$.  Thus the test used in the CLN step is always a difference of two
semipositive slot-zero metrics in the domain of
\cref{cor:uniform-family-r3}; the arbitrary continuous limit
$\psi$ is never inserted into that slot.
Let $C_q$ be the CLN constant for $H_q,L_1,\ldots,L_n$, with $C_q=0$
when $n=0$.  No boundedness of $C_q$ is assumed.  For every $i$, choose
a net of tropical FS metrics $\phi_{i,a_i}\to\phi_i$ uniformly.

Let
\[
 P=(0,1]^{\rm op}\times Q\times\prod_{i=1}^nA_i
\]
with coordinatewise order, and let $D\subset P$ consist of
$d=(\eps,q,\mathbf a)$ such that
\begin{align}
 (1+M)\norm{\psi_q-\psi}_\infty&\leq\eps/2,
\tag{5.1}\label{eq:D-test-r3}\\
 (1+C_q)\sum_{i=1}^n
 \norm{\phi_{i,a_i}-\phi_i}_\infty&\leq\eps/2.
\tag{5.2}\label{eq:D-metric-r3}
\end{align}

The set $D$ is directed.  Given two elements, choose a smaller
$\eps$, then a common later $q$ satisfying~\eqref{eq:D-test-r3}, and,
for that particular $q$, common later $a_i$'s satisfying
\eqref{eq:D-metric-r3}.

Here is the literal cofinality check.  Given an arbitrary ambient index
$p_0=(\eps_0,q_0,\mathbf a_0)\in P$, first choose
$0<\eps\leq\eps_0$.  Uniform convergence of the DFS net gives
$q\geq q_0$ satisfying~\eqref{eq:D-test-r3}; after this $q$ is fixed,
uniform convergence of each FS net gives $a_i\geq a_{0,i}$ satisfying
\eqref{eq:D-metric-r3}.  Hence
$d=(\eps,q,\mathbf a)\in D$ and $d\geq p_0$.  Thus $D$ is genuinely
cofinal in $P$.

For $d\in D$, put
\[
 F_d(x):=\int_{X_x^{\rm an}}\psi_q\,
 \MA_x(\phi_{1,a_1},\ldots,\phi_{n,a_n}),
 \qquad
 F(x):=\int_{X_x^{\rm an}}\psi\,
 \MA_x(\phi_1,\ldots,\phi_n).
\]
Each $F_d$ is continuous by the finite-level hypothesis.  Positivity and
common mass control the test error involving $\psi_q-\psi$, while
\cref{cor:uniform-family-r3}, applied only with the admissible
slot-zero pair $(\alpha_q,\beta_q)$, controls the metric error.  Uniformly
in $x\in\cM(A)$,
\[
 \abs{F_d(x)-F(x)}
 \leq M\norm{\psi_q-\psi}_\infty
 +C_q\sum_i\norm{\phi_{i,a_i}-\phi_i}_\infty
 \leq\eps.
\]
Thus $F_d\to F$ uniformly.  A uniform limit of a net of continuous
real-valued functions is continuous.  The construction neither uses a
sequence nor assumes that the $C_q$ are uniformly bounded.

\section{Exact source and dependency ledger}

\begin{description}[style=nextline]
\item[Pille--Schneider convention.]
Proposition~3.2, printed page~24, gives the two exponent identities reproduced
in~\eqref{eq:PS32-transport}.  Definition~4.8, printed page~36, defines the
family measure by restriction to $X_{\mathcal H(x)}^{\rm an}$.  Their
Theorem~4.11 applies the usual complex current to the unpowered metric
$\phi_t$ from Proposition~3.2.  Their Theorem~2.11 and Proposition~2.21
transport the psh class and its uniform density.
\item[Boucksom--Eriksson normalization.]
Section~1.1 identifies an Archimedean complete valued field with $\R$ or
$\C$ up to normalization.  Example~4.2 identifies its analytic space with
the usual complex space (or real conjugation quotient).  Proposition~8.3
fixes the algebraic mass and real push-forward; Theorem~8.14 supplies the
arbitrary-valued-field Poincare--Lelong compatibility test; Theorem~8.4 gives
continuity.
\item[PN2 and the dimension-zero metric.]
Boucksom--Eriksson Example~8.15, formula~(8.7), identifies the algebraic
norm line, while Ducrot Proposition~6.1.1(3) is used only for its positive
algebraic orientation.  The metric normalization is formula~(8.11) in the
proof of Boucksom--Eriksson Theorem~8.16: in the extracted primary source it
is at lines~3038--3055 (printed page~63).  Example~8.2 and Remark~8.1 are at
extracted lines~2634--2646 (printed page~55), giving respectively the exact
zero-cycle masses and the real push-forward.  The higher-dimensional
recursion and uniqueness are Theorem~8.16, formulas~(8.8)--(8.9).
\item[Weighted Philippon--Gauss.]
The theorem is a frozen exact-byte, production-verified dependency in the
Part~6 fragment directory, in
\path{ch66_weighted_philippon_gauss_statement.tex} and
\path{ch66_weighted_philippon_gauss_proof.tex}.
Those fragments are inputs of the live Chapter~66.  The focused
reconstruction in this repair module is not asserted to be an exact-byte copy
of the verified proof fragment.  That fragment imports the fixed-type
coefficient estimate from D'Andrea--Krick--Sombra and coefficient continuity
from Zhang.
\item[CLN and density.]
Boucksom--Eriksson Lemma~8.6 (extracted lines~2708--2731, printed
pages~56--57) assumes semipositive pairs in every slot, including slot zero.
Pille--Schneider Lemma~4.10 (extracted lines~2063--2070, printed page~37)
likewise assumes FS pairs in every slot at the non-Archimedean endpoint.
DFS and FS density are Pille--Schneider Theorem~2.16 and
Proposition~2.21; the product-net argument uses CLN only on the DFS
decomposition pair, never on the arbitrary continuous limit test.
\end{description}

The convention transport is proved here for the metrics in the
Pille--Schneider psh/DFS closure consumed by the bridge.  The module does not
assert a comparison for a different, independently normalized complex
$dd^c$ convention.  It also does not construct the relative resultant
tensor or prove its comparison with the Chen--Moriwaki fundamental-cycle
resultant; those remain separate P17.1 dependencies.

\begin{thebibliography}{9}
\raggedright
\bibitem{BE}
S.~Boucksom and D.~Eriksson,
\emph{Spaces of norms, determinant of cohomology and Fekete points in
non-Archimedean geometry}, arXiv:1805.01016v3, especially Sections~1.1 and
4.1, Proposition~8.3, Theorem~8.4, Lemma~8.6, Theorem~8.14 and
Theorem~8.16.

\smallskip
\bibitem{CM}
H.~Chen and A.~Moriwaki,
\emph{Arithmetic intersection theory over adelic curves},
arXiv:2103.15646v1, Theorem~3.8.7.  Their Archimedean formula uses the usual
absolute value; PN2 performs the powered transfer.

\smallskip
\bibitem{DKS}
C.~D'Andrea, T.~Krick and M.~Sombra,
\emph{Heights of varieties in multiprojective spaces and arithmetic Nullstellens\"atze},
Ann. Sci. \emph{Ec. Norm. Sup\'er.} 46 (2013), especially Lemma~2.32,
Definition~2.33 and formulas (2.34)--(2.35).

\smallskip
\bibitem{Ducrot}
F.~Ducrot,
\emph{Structures du cube et fibres d'intersection},
arXiv:alg-geom/9712010v1, Proposition~6.1.1.

\smallskip
\bibitem{PS}
L.~Pille-Schneider,
\emph{Global pluripotential theory on hybrid spaces},
J. \emph{Ec. polytech. Math.} 10 (2023), 601--658, especially
Proposition~3.2, Definition~4.8, Theorems~4.3 and 4.11, and Lemma~4.10.

\smallskip
\bibitem{Zhang}
S.~Zhang,
\emph{Heights and reductions of semi-stable varieties},
Compos. Math. 104 (1996), Section~1.5 and Theorem~3.6(1).
\end{thebibliography}

\end{document}
